% ch06_schauder.tex — Th\'eor\`eme de Schauder et Variantes
\chapter{Th\'eor\`eme de Schauder et Variantes}\label{ch:schauder}

\section{Motivation}

Le th\'eor\`eme de Brouwer est un bijou de la topologie en dimension finie, mais il \'echoue spectaculairement en dimension infinie~: la boule unit\'e ferm\'ee d'un espace de Banach de dimension infinie n'est pas compacte, et des applications continues sans point fixe existent. Juliusz Schauder, math\'ematicien polonais de l'\'ecole de Lw\'ow, a surmont\'e cet obstacle en 1930 en rempla\c{c}ant la compacit\'e de l'ensemble par la \emph{compacit\'e de l'op\'erateur}~: si $T$ envoie un convexe ferm\'e born\'e dans un sous-ensemble relativement compact de lui-m\^eme, alors $T$ a un point fixe. Ce r\'esultat est le pilier des m\'ethodes d'existence pour les \'equations diff\'erentielles et int\'egrales non lin\'eaires.

Le th\'eor\`eme de Brouwer (Chapitre~5) ne s'applique qu'en dimension finie.
En dimension infinie, la boule unit\'e ferm\'ee n'est pas compacte, et des
applications continues sans point fixe existent. Schauder a r\'esolu ce
probl\`eme en rempla\c{c}ant la compacit\'e de l'ensemble par celle de l'op\'erateur.

\section{Op\'erateurs compacts}

\begin{definition}[Op\'erateur compact]
Soit $E$ un espace de Banach et $C \subset E$. Un op\'erateur continu
$T : C \to E$ est \emph{compact} si $T(B)$ est relativement compact (c'est-\`a-dire
$\overline{T(B)}$ est compact) pour tout sous-ensemble born\'e $B \subset C$.
\end{definition}

\begin{definition}[Op\'erateur compl\`etement continu]
Un op\'erateur $T : C \to E$ est \emph{compl\`etement continu} s'il est continu et
envoie tout ensemble born\'e sur un ensemble relativement compact.
\end{definition}

\begin{remarque}
Dans un espace de Banach, tout op\'erateur compl\`etement continu est compact. La
r\'eciproque est vraie si le domaine est born\'e.
\end{remarque}

\begin{exemple}[Op\'erateur int\'egral de Fredholm]
L'op\'erateur $(Tu)(t) = \int_0^1 K(t,s) u(s) \, ds$ sur $C([0,1])$ est compact
d\`es que $K \in C([0,1]^2)$. C'est une cons\'equence du th\'eor\`eme d'Arzel\`a--Ascoli.
\end{exemple}

\section{Th\'eor\`eme de Schauder}

\begin{theoreme}[Schauder, 1930]\label{thm:schauder}
Soit $E$ un espace de Banach, $C \subset E$ un sous-ensemble convexe, ferm\'e et
born\'e, et $T : C \to C$ un op\'erateur compact. Alors $T$ poss\`ede au moins un
point fixe.
\end{theoreme}

La preuve repose sur le lemme d'approximation suivant.

\begin{lemme}[Approximation de Schauder]\label{lem:schauder-approx}
Soit $K$ un sous-ensemble compact d'un espace de Banach $E$ et $\varepsilon > 0$.
Il existe un sous-espace de dimension finie $F \subset E$ et une application
continue $P_\varepsilon : K \to \mathrm{co}(K) \cap F$ telle que
$\norm{P_\varepsilon(x) - x} < \varepsilon$ pour tout $x \in K$.
\end{lemme}

\begin{proof}
Par compacit\'e, il existe $x_1, \ldots, x_n \in K$ tels que
$K \subset \bigcup_{i=1}^n B(x_i, \varepsilon)$.
D\'efinissons $\mu_i(x) = \max(0, \varepsilon - \norm{x - x_i})$ et
\[
    P_\varepsilon(x) = \frac{\sum_{i=1}^n \mu_i(x) \, x_i}{\sum_{i=1}^n \mu_i(x)}.
\]
Cette application est bien d\'efinie (le d\'enominateur est positif car $x$ est
\`a distance $< \varepsilon$ d'au moins un $x_i$), continue, et
$P_\varepsilon(x) \in \mathrm{co}(\{x_1, \ldots, x_n\})$. De plus,
\[
    \norm{P_\varepsilon(x) - x} = \norm{\frac{\sum \mu_i(x)(x_i - x)}{\sum \mu_i(x)}}
    \leq \frac{\sum \mu_i(x) \norm{x_i - x}}{\sum \mu_i(x)} < \varepsilon
\]
car $\mu_i(x) > 0$ implique $\norm{x_i - x} < \varepsilon$.
\end{proof}

\begin{proof}[Preuve du Th\'eor\`eme~\ref{thm:schauder}]
Posons $K = \overline{T(C)}$, qui est compact. Pour chaque $n \geq 1$,
soit $P_n = P_{1/n}$ l'approximation de Schauder du Lemme~\ref{lem:schauder-approx}.
L'application $P_n \circ T : C \to \mathrm{co}(K) \cap F_n$ est continue et
\`a valeurs dans un espace de dimension finie.

Posons $C_n = C \cap F_n$, qui est convexe, ferm\'e et born\'e dans $F_n$.
L'application $P_n \circ T|_{C_n} : C_n \to C_n$ est continue (en ajustant le
domaine si n\'ecessaire pour assurer la stabilit\'e). Par le th\'eor\`eme de
Brouwer, il existe $x_n \in C_n$ tel que $P_n(T(x_n)) = x_n$.

On a $\norm{x_n - T(x_n)} = \norm{P_n(T(x_n)) - T(x_n)} < 1/n$. Par compacit\'e
de $T$, la suite $(T(x_n))$ admet une sous-suite convergente~: $T(x_{n_k}) \to y$.
Alors $x_{n_k} \to y$ aussi, et par continuit\'e de $T$~: $T(y) = \lim T(x_{n_k}) = y$.
Puisque $C$ est ferm\'e, $y \in C$.
\end{proof}

\begin{formules}[title=\textbf{Th\'eor\`eme de Schauder}]
$C$ convexe ferm\'e born\'e dans un Banach, $T : C \to C$ compact $\implies$
$\exists\, x^* \in C : T(x^*) = x^*$.
\end{formules}

\section{Version de Schauder--Tychonoff}

\begin{theoreme}[Schauder--Tychonoff]
Soit $E$ un espace localement convexe s\'epar\'e, $C \subset E$ un sous-ensemble
convexe compact non vide, et $T : C \to C$ continue. Alors $T$ poss\`ede un
point fixe.
\end{theoreme}

\begin{remarque}
Ce th\'eor\`eme g\'en\'eralise \`a la fois Brouwer (dimension finie) et Schauder
(espace de Banach). L'hypoth\`ese de convexit\'e locale de l'espace ambiant est
essentielle.
\end{remarque}

\section{Th\'eor\`eme de Leray--Schauder}

\begin{definition}[Degr\'e de Leray--Schauder]
Soit $E$ un espace de Banach, $\Omega \subset E$ un ouvert born\'e, et
$T : \overline{\Omega} \to E$ un op\'erateur compact tel que
$x \neq T(x)$ pour tout $x \in \partial\Omega$. Le \emph{degr\'e de
Leray--Schauder} $\deg_{LS}(\mathrm{Id} - T, \Omega, 0) \in \Z$ est d\'efini
par approximation en dimension finie.
\end{definition}

\begin{theoreme}[Leray--Schauder, 1934]\label{thm:leray-schauder}
Soit $E$ un espace de Banach, $\Omega \subset E$ un ouvert born\'e contenant
$0$, et $T : \overline{\Omega} \to E$ un op\'erateur compact. Supposons que
\[
    x \neq \lambda T(x) \quad \forall x \in \partial\Omega, \; \forall \lambda \in [0, 1].
\]
Alors $T$ poss\`ede un point fixe dans $\Omega$.
\end{theoreme}

\begin{proof}[Id\'ee de la preuve]
L'homotopie $H(\lambda, x) = x - \lambda T(x)$ v\'erifie $H(\lambda, x) \neq 0$
pour $x \in \partial\Omega$ et $\lambda \in [0,1]$. Par invariance homotopique
du degr\'e~:
\[
    \deg_{LS}(\mathrm{Id} - T, \Omega, 0) = \deg_{LS}(\mathrm{Id}, \Omega, 0) = 1.
\]
Puisque le degr\'e est non nul, l'\'equation $x - T(x) = 0$ a une solution.
\end{proof}

\begin{corollaire}[Alternative de Leray--Schauder]
Soit $E$ un espace de Banach et $T : E \to E$ un op\'erateur compl\`etement
continu. Alors, au moins l'une des deux alternatives suivantes est v\'erifi\'ee~:
\begin{enumerate}
    \item L'\'equation $x = T(x)$ a une solution;
    \item L'ensemble $\{x \in E : x = \lambda T(x) \text{ pour un } \lambda \in (0,1)\}$
          est non born\'e.
\end{enumerate}
\end{corollaire}

\section{Mesure de non-compacit\'e et th\'eor\`eme de Darbo}

\begin{definition}[Mesure de non-compacit\'e de Kuratowski]
Soit $A$ un sous-ensemble born\'e d'un espace m\'etrique. La \emph{mesure de
non-compacit\'e de Kuratowski} est
\[
    \alpha(A) = \inf\left\{ \varepsilon > 0 : A \subset \bigcup_{i=1}^n A_i
    \text{ avec } \mathrm{diam}(A_i) \leq \varepsilon \right\}.
\]
\end{definition}

\begin{proposition}[Propri\'et\'es de $\alpha$]
Pour des sous-ensembles born\'es $A, B$ d'un espace de Banach~:
\begin{enumerate}
    \item $\alpha(A) = 0 \Leftrightarrow \overline{A}$ est compact;
    \item $A \subset B \Rightarrow \alpha(A) \leq \alpha(B)$;
    \item $\alpha(A \cup B) = \max(\alpha(A), \alpha(B))$;
    \item $\alpha(\overline{\mathrm{co}}(A)) = \alpha(A)$;
    \item $\alpha(A + B) \leq \alpha(A) + \alpha(B)$;
    \item $\alpha(\lambda A) = \abs{\lambda} \alpha(A)$.
\end{enumerate}
\end{proposition}

\begin{definition}[Application condensante]
Un op\'erateur continu $T : C \to C$ (o\`u $C$ est convexe ferm\'e born\'e) est
\emph{$\alpha$-condensant} s'il existe $k \in [0, 1)$ tel que
$\alpha(T(A)) \leq k \, \alpha(A)$ pour tout $A \subset C$ born\'e.
\end{definition}

\begin{theoreme}[Darbo, 1955]\label{thm:darbo}
Soit $E$ un espace de Banach, $C \subset E$ convexe ferm\'e born\'e non vide,
et $T : C \to C$ un op\'erateur $\alpha$-condensant. Alors $T$ poss\`ede un
point fixe.
\end{theoreme}

\begin{proof}
D\'efinissons par r\'ecurrence $C_0 = C$ et $C_{n+1} = \overline{\mathrm{co}}(T(C_n))$.
La suite $(C_n)$ est d\'ecroissante ($C_{n+1} \subset C_n$) car $T(C_n) \subset C_n$
implique $C_{n+1} = \overline{\mathrm{co}}(T(C_n)) \subset \overline{\mathrm{co}}(C_n) = C_n$.

Posons $\alpha_n = \alpha(C_n)$. Alors
$\alpha_{n+1} = \alpha(\overline{\mathrm{co}}(T(C_n))) = \alpha(T(C_n)) \leq k \alpha_n$.
Donc $\alpha_n \leq k^n \alpha_0 \to 0$.

Posons $C_\infty = \bigcap_{n \geq 0} C_n$. Comme intersection d\'ecroissante de
ferm\'es dans un espace complet avec $\alpha(C_n) \to 0$, l'ensemble $C_\infty$
est compact non vide (par Cantor g\'en\'eralis\'e). De plus, $C_\infty$ est convexe.
On a $T(C_\infty) \subset T(C_n) \subset C_{n+1}$ pour tout $n$, donc
$T(C_\infty) \subset C_\infty$.

Par le th\'eor\`eme de Schauder appliqu\'e \`a $T : C_\infty \to C_\infty$ (continu
sur un convexe compact), $T$ a un point fixe.
\end{proof}

\begin{remarque}
Le th\'eor\`eme de Schauder est un cas particulier de celui de Darbo~: si $T$ est
compact, alors $\alpha(T(A)) = 0$ pour tout born\'e $A$, donc $T$ est
$\alpha$-condensant avec $k = 0$.
\end{remarque}

\section{Th\'eor\`eme de Sadovskii}

\begin{theoreme}[Sadovskii, 1967]
Soit $C$ un sous-ensemble convexe, ferm\'e et born\'e d'un espace de Banach, et
$T : C \to C$ continue telle que $\alpha(T(A)) < \alpha(A)$ pour tout
$A \subset C$ avec $\alpha(A) > 0$. Alors $T$ a un point fixe.
\end{theoreme}

\begin{remarque}
Sadovskii g\'en\'eralise Darbo en rempla\c{c}ant la condition
$\alpha(T(A)) \leq k\alpha(A)$ ($k < 1$) par $\alpha(T(A)) < \alpha(A)$.
La preuve utilise un argument de Zorn.
\end{remarque}

\section{Th\'eor\`eme de Krasnoselskii}

\begin{theoreme}[Krasnoselskii, 1958]\label{thm:krasnoselskii}
Soit $C$ un sous-ensemble convexe, ferm\'e et born\'e d'un espace de Banach $E$.
Soient $A : C \to E$ un op\'erateur compact et $B : C \to E$ une contraction
de constante $k < 1$. Si $A(x) + B(y) \in C$ pour tous $x, y \in C$, alors
l'\'equation $x = A(x) + B(x)$ a une solution dans $C$.
\end{theoreme}

\begin{proof}
Pour chaque $x \in C$, l'application $y \mapsto A(x) + B(y)$ est une contraction
de $C$ dans $C$ (par hypoth\`ese). Par Banach, elle a un unique point fixe
$y(x)$~: $y(x) = A(x) + B(y(x))$. L'application $x \mapsto y(x)$ est continue
(par d\'ependance continue du point fixe par rapport au param\`etre) et compacte
(car $y(x) = A(x) + B(y(x))$ et $A$ est compact). Par Schauder, $y(\cdot)$ a
un point fixe $x^* = y(x^*)$, et donc $x^* = A(x^*) + B(x^*)$.
\end{proof}

\begin{remarque}
Le th\'eor\`eme de Krasnoselskii est particuli\`erement utile pour les \'equations
int\'egrales de la forme $x = Ax + Bx$ o\`u $A$ est un op\'erateur int\'egral
(compact) et $B$ est une contraction.
\end{remarque}

\section{Applications}

\begin{exemple}[\'Equation int\'egrale non lin\'eaire]
Consid\'erons l'\'equation
\[
    u(t) = g(t) + \int_0^1 K(t,s) \, h(s, u(s)) \, ds, \quad t \in [0,1]
\]
o\`u $g \in C([0,1])$, $K \in C([0,1]^2)$ et $h : [0,1] \times \R \to \R$ est
continue et born\'ee. L'op\'erateur
$(Tu)(t) = g(t) + \int_0^1 K(t,s) h(s, u(s)) \, ds$ est compact sur la boule
$\overline{B}(0, R) \subset C([0,1])$ pour $R$ assez grand, et
$T(\overline{B}(0,R)) \subset \overline{B}(0,R)$ si $R \geq \norm{g}_\infty +
\norm{K}_\infty \norm{h}_\infty$. Par Schauder, l'\'equation a une solution.
\end{exemple}

\section{Exercices}

\begin{exercice}
Montrer que l'op\'erateur de Volterra $(Tu)(t) = \int_0^t K(t,s) u(s) \, ds$
est compact sur $C([0,1])$ lorsque $K \in C([0,1]^2)$.
\end{exercice}

\begin{exercice}
V\'erifier que $\alpha(\overline{B}(0, r)) = 2r$ dans un espace de Banach de
dimension infinie, et $\alpha(\overline{B}(0, r)) = 0$ en dimension finie.
\end{exercice}

\begin{exercice}
Utiliser le th\'eor\`eme de Schauder pour prouver l'existence d'une solution au
probl\`eme de Cauchy--Peano~: $y' = f(t,y)$, $y(0) = y_0$, o\`u $f$ est continue
(mais pas n\'ecessairement lipschitzienne).
\end{exercice}

\begin{exercice}
Montrer que le th\'eor\`eme de Darbo implique celui de Schauder.
\end{exercice}

\begin{exercice}
Appliquer le th\'eor\`eme de Krasnoselskii pour montrer l'existence d'une solution
\`a l'\'equation $u(t) = \frac{1}{3} u(t) + \int_0^1 K(t,s) g(u(s)) \, ds$ dans
$C([0,1])$, sous des hypoth\`eses convenables sur $K$ et $g$.
\end{exercice}

\begin{exercice}
Montrer que dans un espace de Banach r\'eflexif, la boule unit\'e ferm\'ee est
faiblement compacte. En d\'eduire une version du th\'eor\`eme de Schauder pour la
topologie faible.
\end{exercice}
